\documentclass[10pt,a4paper]{amsart}
\usepackage{amsmath,amssymb,amsthm,graphicx}
\usepackage{verbatim}
\newenvironment{reference}{\comment}{\endcomment}
%\newenvironment{reference}{}{}
\newenvironment{slogan}{\comment}{\endcomment}
\newenvironment{history}{\comment}{\endcomment}

% For commutative diagrams we use Xy-pic
\usepackage[all]{xy}

% We use 2cell for 2-commutative diagrams.
\xyoption{2cell}
% Theorem environments.
%
\theoremstyle{plain}
\newtheorem{theorem}[subsection]{Theorem}
\newtheorem{proposition}[subsection]{Proposition}
\newtheorem{lemma}[subsection]{Lemma}

\theoremstyle{definition}
\newtheorem{definition}[subsection]{Definition}
\newtheorem{example}[subsection]{Example}
\newtheorem{exercise}[subsection]{Exercise}
\newtheorem{situation}[subsection]{Situation}

\theoremstyle{remark}
\newtheorem{remark}[subsection]{Remark}
\newtheorem{remarks}[subsection]{Remarks}

\numberwithin{equation}{subsection}

% Macros
%
\def\lim{\mathop{\mathrm{lim}}\nolimits}
\def\colim{\mathop{\mathrm{colim}}\nolimits}
\def\Spec{\mathop{\mathrm{Spec}}}
\def\Hom{\mathop{\mathrm{Hom}}\nolimits}
\def\Ext{\mathop{\mathrm{Ext}}\nolimits}
\def\SheafHom{\mathop{\mathcal{H}\!\mathit{om}}\nolimits}
\def\SheafExt{\mathop{\mathcal{E}\!\mathit{xt}}\nolimits}
\def\Sch{\mathit{Sch}}
\def\Mor{\mathop{\mathrm{Mor}}\nolimits}
\def\Ob{\mathop{\mathrm{Ob}}\nolimits}
\def\Sh{\mathop{\mathit{Sh}}\nolimits}
\def\NL{\mathop{N\!L}\nolimits}
\def\CH{\mathop{\mathrm{CH}}\nolimits}
\def\proetale{{pro\text{-}\acute{e}tale}}
\def\etale{{\acute{e}tale}}
\def\QCoh{\mathit{QCoh}}
\def\Ker{\mathop{\mathrm{Ker}}}
\def\Im{\mathop{\mathrm{Im}}}
\def\Coker{\mathop{\mathrm{Coker}}}
\def\Coim{\mathop{\mathrm{Coim}}}

% Boxtimes
%
\DeclareMathSymbol{\boxtimes}{\mathbin}{AMSa}{"02}

%
% Macros for moduli stacks/spaces
%
\def\QCohstack{\mathcal{QC}\!\mathit{oh}}
\def\Cohstack{\mathcal{C}\!\mathit{oh}}
\def\Spacesstack{\mathcal{S}\!\mathit{paces}}
\def\Quotfunctor{\mathrm{Quot}}
\def\Hilbfunctor{\mathrm{Hilb}}
\def\Curvesstack{\mathcal{C}\!\mathit{urves}}
\def\Polarizedstack{\mathcal{P}\!\mathit{olarized}}
\def\Complexesstack{\mathcal{C}\!\mathit{omplexes}}
% \Pic is the operator that assigns to X its picard group, usage \Pic(X)
% \Picardstack_{X/B} denotes the Picard stack of X over B
% \Picardfunctor_{X/B} denotes the Picard functor of X over B
\def\Pic{\mathop{\mathrm{Pic}}\nolimits}
\def\Picardstack{\mathcal{P}\!\mathit{ic}}
\def\Picardfunctor{\mathrm{Pic}}
\def\Deformationcategory{\mathcal{D}\!\mathit{ef}}
\usepackage{cleveref}
\crefname{originallemma}{original lemma}{original lemmas}
\Crefname{originallemma}{Original lemma}{Original lemmas}
\makeatletter
\newlabel{algebra-lemma-dominate}{{00IA}{source}{}{sourcepointer-1}{}}
\newlabel{algebra-lemma-dominate@cref}{{[originallemma][0][]00IA}{[1][0][]source}{}{}{}}
\makeatother
\crefname{originallemma}{original lemma}{original lemmas}
\Crefname{originallemma}{Original lemma}{Original lemmas}
\makeatletter
\newlabel{algebra-lemma-image-stable-specialization-closed}{{00HY}{source}{}{sourcepointer-2}{}}
\newlabel{algebra-lemma-image-stable-specialization-closed@cref}{{[originallemma][0][]00HY}{[1][0][]source}{}{}{}}
\makeatother
\crefname{originallemma}{original lemma}{original lemmas}
\Crefname{originallemma}{Original lemma}{Original lemmas}
\makeatletter
\newlabel{lemma-base-change-quasi-compact}{{03HF}{source}{}{sourcepointer-3}{}}
\newlabel{lemma-base-change-quasi-compact@cref}{{[originallemma][0][]03HF}{[1][0][]source}{}{}{}}
\makeatother
\crefname{originallemma}{original lemma}{original lemmas}
\Crefname{originallemma}{Original lemma}{Original lemmas}
\makeatletter
\newlabel{lemma-base-change-valuative-criteria}{{03IY}{source}{}{sourcepointer-4}{}}
\newlabel{lemma-base-change-valuative-criteria@cref}{{[originallemma][0][]03IY}{[1][0][]source}{}{}{}}
\makeatother
\crefname{originallemma}{original lemma}{original lemmas}
\Crefname{originallemma}{Original lemma}{Original lemmas}
\makeatletter
\newlabel{lemma-universally-closed-local}{{03IT}{source}{}{sourcepointer-5}{}}
\newlabel{lemma-universally-closed-local@cref}{{[originallemma][0][]03IT}{[1][0][]source}{}{}{}}
\makeatother
\crefname{originalsection}{original section}{original sections}
\Crefname{originalsection}{Original section}{Original sections}
\makeatletter
\newlabel{schemes-section-points}{{01J5}{source}{}{sourcepointer-6}{}}
\newlabel{schemes-section-points@cref}{{[originalsection][0][]01J5}{[1][0][]source}{}{}{}}
\makeatother
\crefname{originallemma}{original lemma}{original lemmas}
\Crefname{originallemma}{Original lemma}{Original lemmas}
\makeatletter
\newlabel{spaces-properties-lemma-quasi-compact-affine-cover}{{03H6}{source}{}{sourcepointer-7}{}}
\newlabel{spaces-properties-lemma-quasi-compact-affine-cover@cref}{{[originallemma][0][]03H6}{[1][0][]source}{}{}{}}
\makeatother
\title{A quasi-compact morphism of algebraic spaces satisfying the existence part of the valuative criterion is universally closed}
\author{The Stacks Project Authors}
\date{stacks-project:spaces-morphisms.tex:local-snapshot-20261012; GNU FDL 1.2 or later}
\begin{document}
\maketitle
\setcounter{subsection}{1}
\noindent\textit{Editorial note.} All delivered mathematical text is the author's own, copied from the pinned Stacks Project chapter file spaces-morphisms.tex; no statement or proof step was rewritten, reordered or removed.\par\medskip
\noindent\textit{Editorial note.} The proof is the author's own valuative argument on points; the closedness conclusion is obtained by specialization, not by a citation of the scheme case.\par\medskip
\noindent\textit{Editorial note.} The seven references of the counted spans are recorded as non-load-bearing navigation pointers; the argument that consumes them stays inside the counted proof span.\par\medskip
\noindent\textit{Editorial note.} Preamble: the three native excerpt spans of sources/stacks-spaces-morphisms/preamble.tex (lines 8-13, 14-18 and 154-217). The chapter file itself carries no preamble text.\par\medskip
\noindent\textit{Allowed standard prerequisites (editorial).}\begin{itemize}
\item Universally closed morphisms are detected on affine schemes mapping into the base. Lemma lemma-universally-closed-local, used at proof lines 8813-8815. Reduces to showing that |T x\_Y X| -> |T| is closed for affine T over S.
\item Stability of quasi-compactness and of the valuative criterion under base change. Lemmas lemma-base-change-quasi-compact and lemma-base-change-valuative-criteria, used at proof lines 8810-8813. Lets the hypotheses be assumed for the base changed morphism.
\item Closedness of a subset of an affine scheme stable under specialization. Algebra, Lemma algebra-lemma-image-stable-specialization-closed, used at proof lines 8825-8828. Reduces the closedness of f(T) to stability under specialization.
\item Existence of a valuation ring dominating a local ring inside a field. Algebra, Lemma algebra-lemma-dominate, used at proof lines 8844-8847. Produces the valuation ring A in K dominating the image of O\_\{Y,y\}.
\item An affine etale cover of a quasi-compact algebraic space. Properties of Spaces, Lemma spaces-properties-lemma-quasi-compact-affine-cover, used at proof lines 8819-8823. Turns the closed subset T into the preimage of a closed subscheme of an affine scheme.
\end{itemize}\medskip
\ifdefined\hypertarget\hypertarget{sourcepointer-1}{}\fi
\noindent\textit{Source reference.} Original lemma 00IA: focused-seeds/AGP/sources/stacks-snapshot-20261011/tags.txt, source lines 672--672. Navigation pointer into the Stacks Project; the standard fact it supplies is applied inside the counted proof span and the argument resting on it stays there.\par\medskip
\ifdefined\hypertarget\hypertarget{sourcepointer-2}{}\fi
\noindent\textit{Source reference.} Original lemma 00HY: focused-seeds/AGP/sources/stacks-snapshot-20261011/tags.txt, source lines 660--660. Navigation pointer into the Stacks Project; the standard fact it supplies is applied inside the counted proof span and the argument resting on it stays there.\par\medskip
\ifdefined\hypertarget\hypertarget{sourcepointer-3}{}\fi
\noindent\textit{Source reference.} Original lemma 03HF: focused-seeds/AGP/sources/stacks-snapshot-20261011/tags.txt, source lines 4565--4565. Navigation pointer into the Stacks Project; the standard fact it supplies is applied inside the counted proof span and the argument resting on it stays there.\par\medskip
\ifdefined\hypertarget\hypertarget{sourcepointer-4}{}\fi
\noindent\textit{Source reference.} Original lemma 03IY: focused-seeds/AGP/sources/stacks-snapshot-20261011/tags.txt, source lines 4620--4620. Navigation pointer into the Stacks Project; the standard fact it supplies is applied inside the counted proof span and the argument resting on it stays there.\par\medskip
\ifdefined\hypertarget\hypertarget{sourcepointer-5}{}\fi
\noindent\textit{Source reference.} Original lemma 03IT: focused-seeds/AGP/sources/stacks-snapshot-20261011/tags.txt, source lines 4615--4615. Navigation pointer into the Stacks Project; the standard fact it supplies is applied inside the counted proof span and the argument resting on it stays there.\par\medskip
\ifdefined\hypertarget\hypertarget{sourcepointer-6}{}\fi
\noindent\textit{Source reference.} Original section 01J5: focused-seeds/AGP/sources/stacks-snapshot-20261011/tags.txt, source lines 2002--2002. Navigation pointer into the Stacks Project; the standard fact it supplies is applied inside the counted proof span and the argument resting on it stays there.\par\medskip
\ifdefined\hypertarget\hypertarget{sourcepointer-7}{}\fi
\noindent\textit{Source reference.} Original lemma 03H6: focused-seeds/AGP/sources/stacks-snapshot-20261011/tags.txt, source lines 4555--4555. Navigation pointer into the Stacks Project; the standard fact it supplies is applied inside the counted proof span and the argument resting on it stays there.\par\medskip
\begin{lemma}
\label{lemma-quasi-compact-existence-universally-closed}
Let $S$ be a scheme.
Let $f : X \to Y$ be a morphism of algebraic spaces over $S$.
Assume
\begin{enumerate}
\item $f$ is quasi-compact, and
\item $f$ satisfies the existence part of the valuative criterion.
\end{enumerate}
Then $f$ is universally closed.
\end{lemma}

\begin{proof}
By Lemmas \ref{lemma-base-change-quasi-compact}
and \ref{lemma-base-change-valuative-criteria}
properties (1) and (2) are preserved under
any base change. By Lemma \ref{lemma-universally-closed-local}
we only have to show that $|T \times_Y X| \to |T|$ is closed,
whenever $T$ is an affine scheme over $S$ mapping into $Y$. Hence it
suffices to prove: If $Y$ is an affine scheme, $f : X \to Y$ is quasi-compact
and satisfies the existence part of the valuative criterion, then
$f : |X| \to |Y|$ is closed. In this situation $X$ is a quasi-compact
algebraic space. By
Properties of Spaces,
Lemma \ref{spaces-properties-lemma-quasi-compact-affine-cover}
there exists an affine scheme $U$ and a surjective \'etale morphism
$\varphi : U \to X$. Let $T \subset |X|$ closed. The inverse image
$\varphi^{-1}(T) \subset U$ is closed, and hence is the set of points
of an affine closed subscheme $Z \subset U$. Thus, by
Algebra, Lemma \ref{algebra-lemma-image-stable-specialization-closed}
we see that $f(T) = f(\varphi(|Z|)) \subset |Y|$ is closed if it is
closed under specialization.

\medskip\noindent
Let $y' \leadsto y$ be a specialization in $Y$ with $y' \in f(T)$.
Choose a point $x' \in T \subset |X|$ mapping to $y'$ under $f$.
We may represent $x'$ by a morphism $\Spec(K) \to X$
for some field $K$. Thus we have the following diagram
$$
\xymatrix{
\Spec(K) \ar[r]_-{x'} \ar[d] & X \ar[d]^f \\
\Spec(\mathcal{O}_{Y, y}) \ar[r] & Y,
}
$$
see
Schemes, Section \ref{schemes-section-points}
for the existence of the left vertical map.
Choose a valuation ring $A \subset K$ dominating the image of
the ring map $\mathcal{O}_{Y, y} \to K$ (this is possible since
the image is a local ring and not a field as $y' \not = y$, see
Algebra, Lemma \ref{algebra-lemma-dominate}).
By assumption there exists a field extension $K'/K$ and a
valuation ring $A' \subset K'$ dominating $A$, and a morphism
$\Spec(A') \to X$ fitting into the commutative diagram.
Since $A'$ dominates $A$, and $A$ dominates $\mathcal{O}_{Y, y}$
we see that the closed point of $\Spec(A')$ maps to
a point $x \in X$ with $f(x) = y$ which is a specialization of $x'$.
Hence $x \in T$ as $T$ is closed, and hence $y \in f(T)$ as desired.
\end{proof}

\end{document}
